\setcounter{figure}{11}
\begin{figure}[t!]
\centering
\includegraphics[width=\textwidth]{fig_qif_timing_twotone.pdf}
\caption{\textbf{Beat-locked modulation and mean-rate change in the QIF network.} \emph{(a)} Field-free power
spectrum of $r_E(t)$, normalized: the spiking network ($N=8000$) matches the exact NMM2
mean field in gamma frequency ($\sim$55~Hz); a normalized spectrum does not test the mean
rate. \emph{(b)} Intermodulation above the
Hopf (TI off vs.\ on, log power): the nonlinearity mixes the imposed $\Delta f=42$~Hz beat
with the persisting gamma near $53$--$54$~Hz, raising the direct $\Delta f$ line and the sum sideband
near $95$~Hz, both absent with TI off. These sidebands establish nonlinear mixing;
they do not uniquely identify gamma-amplitude gating. \emph{(c)} The $\Delta f$ modulation depth
$A_\Omega/\langle r_E\rangle$, off vs.\ on: below $0.003$ rising to $0.29$
($\bar\eta=0.85$) and $0.23$ ($\bar\eta=1.5$), while the mean rate rises by about
$8$--$11\%$. The labels identify field-free states; both field-on responses contain
gamma activity beyond the beat component. Spectra use the same post-settling windows
as the rate measurements and are normalized to each spectrum's maximum. Modulation depth is
reported rather than the off$\to$on lock-in ratio, whose denominator is the field-free
floor.}
\label{fig:qif_timing}
\end{figure}
